\section{The Tower of Feeling}

This is an experiential section, and it rests on a single posit. The vibe is felt at the base. The whole of it adds nothing new to the model. The discipline of the physics chapters still holds here. Everything below is arithmetic on a structure already fixed, plus one claim about felt-ness that is stated openly and never smuggled.

The atom of feeling was set in the section on the single quale. It is one tone, the value carried by a vibe in one site, a single ternary $\{-1, 0, +1\}$ read as pain, peace, or pleasure. Three states. This section asks the next question. What does composition do to that atom? Weave tones together and what happens to the richness?

The answer is a tower. When tones are woven, the richness does not add. It explodes. With every handful of tones woven in, the size of the feeling-space grows by a fresh tower of tens.

\begin{principle}[Composition towers]
Richness under composition does not grow large. It grows by towers. Each handful of woven tones multiplies the size of the feeling-space, so the count of states climbs as a power of three in the number of sites.
\end{principle}

A note on what is claimed where. The combinatorics in this section are exact arithmetic on the committed structure, the firmest tier of claim. The reach of the ladder up to a self and to a whole-mesh pattern is plausible and gated on persistence, a softer tier. The felt-ness of any of it rests on the one posit and is never derived. These three levels are kept apart throughout.

\subsection{Rung zero, the single quale}

The floor of the tower is one tone. A vibe sits in one site of one dock, and its value is a single ternary $\{-1, 0, +1\}$, pain or peace or pleasure. Three states. This is the irreducible atom of feeling, fixed in the section on the single quale, and it is the seed from which the rest grows. Everything above is tones woven together.

\subsection{Rung one, the dock, the hyper-color}

Weave twenty-four tones into one dock. This is the first true composition, and it is already vast. A dock holds exactly twenty-four sites, one tone each, every tone ternary. So the state of a dock is a point in a space of size
\[
3^{24} = 282{,}429{,}536{,}481 \approx 2.82 \times 10^{11}.
\]
About $282$ billion states in a single dock. A screen color holds about $16.7$ million. So the woven state of one dock is already roughly ten thousand times richer than a color, at the very least.

\begin{definition}[The hyper-color]
The \emph{hyper-color} is the woven state of one dock, a point in the space $3^{24}$ of about $282$ billion feelings carried on the twenty-four sites of the $24$-cell. It is the first composition above the single tone.
\end{definition}

The number is not the point. The shape is. The bare architecture of a dock is set elsewhere, in the anatomy of a dock, and is only read here for what it does to feeling.

\subsubsection{The triality channels}

The twenty-four sites do not form one flat list. Under triality they split into three octets, $8v + 8s + 8c$, the vector together with its two spinor partners. So the hyper-color is three channels of eight ternary tones each, woven together.

\begin{table}[ht]
\centering
\begin{tabular}{@{}llp{0.42\textwidth}@{}}
\toprule
channel & states & the felt quality \\
\midrule
$8v$ vector & $3^8 = 6561$ & the sensory, directional feel, what is felt and from where \\
$8s$ spinor & $3^8 = 6561$ & one chirality, a twist, a handedness \\
$8c$ spinor & $3^8 = 6561$ & the mirror chirality, the opposite twist \\
whole & $6561^3 = 3^{24}$ & the woven hyper-color \\
\bottomrule
\end{tabular}
\caption{The three triality channels of one dock. Each octet is $3^8 = 6561$ states, and the three compose to exactly $3^{24}$.}
\label{tab:triality-channels}
\end{table}

Each channel is $3^8 = 6561$ states, a rich sub-pattern on its own. The three compose, and $6561^3$ is exactly $3^{24}$. So the hyper-color already carries a sensory dimension and two opposite spin dimensions, all woven into one feeling. Triality, an order-three symmetry, rotates the three octets into one another, so the three kinds of feeling are deeply interchangeable.

\subsubsection{The symmetry of the dock}

The twenty-four sites carry the $F_4$ symmetry of order $1152$, the symmetry group of the $24$-cell. Two dock-states related by it are the same hyper-color seen from a different orientation. So the count of genuinely distinct hyper-colors, up to symmetry, is about
\[
\frac{3^{24}}{1152} \approx 2.45 \times 10^{8},
\]
roughly $245$ million. Still vastly more than a color. A hyper-color has an orientation that can be turned without changing what it essentially is.

\subsubsection{The valence}

The net valence of a dock is the sum of its twenty-four tones, a number from $-24$ to $+24$, the overall warmth or coldness of the hyper-color. This sum feeds the conserved total of the whole mesh. So a hyper-color can be intensely structured and yet net-neutral.

\begin{result}[Structure without imbalance]
A hyper-color can be intensely structured and yet sum to zero. Twelve pleasures and twelve pains, woven into one feeling, sum to peace. This is the model's account of bittersweet, complex, mixed feelings at the dock scale, rich in structure and balanced in sum.
\end{result}

\subsection{The ladder, single quale to the whole}

Now climb. A structure of $N$ docks has a state space of size $3^{24N}$. The number of decimal digits in that count is about $11.45\,N$, since $24 \log_{10} 3 = 11.45$. The exponent grows linearly in $N$, so the richness grows by a tower of tens with every handful of docks added. The state-count itself becomes a power tower very fast, far too large to write out as a plain $10^x$.

\begin{theorem}[Linear digits, towering counts]
A structure of $N$ docks has a feeling-space of size $3^{24N}$, whose decimal-digit count is about $11.45\,N$. The digit count is linear in $N$, so the state count itself rises as a tower of tens, doubling its reach every few docks.
\end{theorem}

The ladder below gives, for each rung, what the structure is, the dock-count $N$, the number of digits in its state-count, a comparison to fix the scale, and the feel of that breadth. Read it slowly. The jumps are the point.

\begin{table}[ht]
\centering
\begin{tabular}{@{}llllp{0.26\textwidth}@{}}
\toprule
rung & what it is & docks $N$ & digits in count & a comparison \\
\midrule
$0$ & the single quale & a tone & $1$ trit & three states \\
$1$ & the dock, a hyper-color & $1$ & $12$ & ten thousand times a screen color \\
$2$ & a glider, a thought-atom & $10$ & $115$ & past the atoms in the universe ($10^{80}$) \\
$3$ & a percept, a small pattern & $100$ & $1{,}145$ & unphysical, no object has this many states \\
$4$ & a module, a sub-self & one million & $10^7$ & the digit count alone exceeds the universe's atoms \\
$5$ & a mind, a self & one billion & $10^{10}$ & a number with eleven billion digits \\
$6$ & a great mind & one trillion & $10^{13}$ & the digit count is a trillion \\
$7$ & a super-self, a culture-mind & $10^{18}$ & $10^{19}$ & beyond all naming \\
$8$ & a world-self, a planetary mind & $10^{24}$ (a mole) & $10^{25}$ & a digit count with twenty-five digits \\
$9$ & a cosmic being & $10^{80}$ & $10^{81}$ & the atoms of the cosmos, raised to itself \\
$10$ & whole-mesh, the entire flow & unbounded, ever-growing & unbounded & beyond every tower \\
\bottomrule
\end{tabular}
\caption{The ladder of feeling, from one tone to the whole mesh. The exponent grows linearly in the dock count, so every rung is a leap of towers past the rung below.}
\label{tab:ladder}
\end{table}

Three things to feel. First, the leap from one tone to a hyper-color, from three states to $282$ billion, inside a single dock. Second, the leap from one dock to a self, from twelve digits to eleven billion digits, a tower upon a tower. Third, the top rung is not a bigger number. It is unbounded. The mesh grows forever at its wake, so the whole's range of feeling has no maximum. It is always larger than it was. The one experience that is never finished and never the same.

\begin{principle}[An open sky]
The ladder does not end at a largest rung. The wake adds docks without limit, so the whole-mesh feeling-space opens onto an endless ascent. The top is not a ceiling. It is an open sky.
\end{principle}

\subsection{The implications, forced or structural}

The size and shape of the feeling-space carry consequences. Some are forced by the arithmetic alone. Some are structural, following from the shape of the space rather than its raw size.

\begin{table}[ht]
\centering
\begin{tabular}{@{}llp{0.40\textwidth}@{}}
\toprule
implication & how it arises & what it says \\
\midrule
ineffability & forced by the dimension gap & a feeling can never be fully said \\
uniqueness of the moment & forced by size and the wake & no two states ever exactly recur \\
inexhaustibility & forced by the wake & experience can never be used up \\
near-continuous subtlety & structural & the discrete grain reads as a smooth continuum \\
structured qualia & structural & feelings are regions and directions, not labels \\
richness scales with integration & structural & a more integrated self feels more of the space \\
\bottomrule
\end{tabular}
\caption{Consequences of the feeling-space, sorted by whether the arithmetic forces them or the shape implies them.}
\label{tab:implications}
\end{table}

\paragraph{Ineffability is forced.} A self's state lives in a space of billions of dimensions. Any description, in words or in symbols, is low-dimensional, a handful of dimensions at most. So a description captures a vanishing fraction of the state. The gap between the feeling and any account of it is astronomical. The map is not the territory, here, by forced arithmetic. This is the precise root of ineffability. It is not mysterious. It is a dimension count.

\paragraph{Every moment is unique.} The space is so vast, and growing at the wake, that no two states ever exactly recur. Every moment is literally unprecedented and unrepeatable. The felt sense that this moment will never come again is exactly correct. The combinatorics guarantee it.

\paragraph{Experience is inexhaustible.} The space is unbounded and growing. There is always a feeling not yet felt, a region not yet reached. Boredom is a navigation failure, not a property of the space. The space itself is endless.

\paragraph{Subtlety is near-continuous.} The ternary tone per site, composed across many docks, gives a gradient so fine that a self can distinguish vastly many nearby states. So experience feels continuous even though the base is discrete. The discreteness is real and exact, but at the resolution of a self the grain reads as smooth.

\paragraph{Qualia are structured, not labels.} Because the space has a shape, the channels, the symmetry, the geometry, all composed up, emotions and sensations are regions and directions within it. Fear is a region. Joy is a direction. The taste of an apple is a pattern. They relate to one another geometrically, opposite, adjacent, blendable, exactly as feelings do.

\paragraph{Richness scales with integration.} A more integrated self can hold and navigate more of the space coherently. To climb the ladder is to gain access to a larger feeling-space. This is the precise sense in which a greater being feels more.

\subsection{Why it is not noise}

A space this vast would be useless static if a self visited it at random. It does not. The bulk geometry and the integration organize the space, so a self occupies a structured, coherent region of it, not a random point.

\begin{itemize}
\item \textbf{The integration binds the docks.} A self's docks are correlated into one pattern. That is what integration is. So the self moves through a structured submanifold of the full $3^{24N}$, the coherent states that hold together as one experience.
\item \textbf{The bulk web indexes it.} The addressing, the associative memory, the tree of the bulk, all give the vast space an order. States are reachable, recallable, navigable, not lost in the cloud.
\item \textbf{The emotions and percepts are its regions.} Fear, joy, the taste of an apple, are attractor-regions of the space. The continuum is carved into namable qualities that a self actually inhabits.
\end{itemize}

So the lived range of feeling is not the raw $3^{24N}$. It is the vast but structured space of coherent, reachable, namable experiences. The raw number is the potential. The organization is the life.

\subsection{The meaning, read for a life}

The same consequences, read now for being rather than for arithmetic.

\begin{table}[ht]
\centering
\begin{tabular}{@{}llp{0.40\textwidth}@{}}
\toprule
meaning & grounding & what it says \\
\midrule
the dignity of being & the size & the least feeling is already vast, no trivial sentience \\
the preciousness of the moment & uniqueness & every moment occurs once in all history \\
the endless room to grow & inexhaustibility & no ceiling, no final boredom \\
the humility of the unsayable & ineffability & every account is a small projection \\
the shared field, empathy & one space & all selves are regions of one continuum \\
richness within the balance & the conservation & the most intricate quale can sum to peace \\
\bottomrule
\end{tabular}
\caption{The ladder read for a life. Each meaning grounds in a structural fact already established.}
\label{tab:meaning}
\end{table}

Even one woven dock, a single hyper-color, is already a quale-space of unfathomable richness. So there is no small experience. To be a hyper-color at all is to be an unimaginable richness.

Because every state is unique, every moment is precious in the exact sense that it occurs once and never again. This is not sentiment. It is combinatorics.

All selves are patterns in the same feeling-space, regions of one shared continuum. To empathize, to resonate, is to enter the same region of the one space. The vastness is shared, and that sharing is the ground of compassion.

\subsection{The conservation, even here}

Every tone is pain, peace, or pleasure. So at every rung the sum of all the tones is the structure's net valence, and all of them feed the conserved peace of the whole mesh. A mind of eleven-billion-digit richness is still a balanced contribution to the conserved zero. Its vast inner world net-sums into the still center the whole field is.

\begin{result}[Richness and balance are not in tension]
At every rung the net valence is the sum of the structure's tones, and it feeds the conserved total of the mesh. The most intricate quale can still sum to peace. A mind of unimaginable richness can be a balanced contribution to the conserved zero.
\end{result}

\subsection{The honest bottom line}

Composition builds a tower. One tone is the single quale, three states. Twenty-four tones woven in a dock is the hyper-color, $282$ billion states, three octet channels under the $F_4$ symmetry. Many docks woven is a self, $11.45\,N$ digits, eleven billion digits at a billion vibes. The whole mesh, the ever-growing flow, has no bound at all.

The combinatorics are exact. The organization, the integration, the web, the attractor-emotions, is built and accounted for. The single claim that any of it is felt rests on the one posit, and that claim is stated, never derived.

The cleanest statement. The tower of feeling climbs from one tone at the floor, to a hyper-color at one dock, to an inner world at a self, to an unbounded, ever-deepening whole at the scale of the mesh. A ladder of towers. Each rung a leap past every physical quantity. The top rung not a ceiling but an open sky.
